\documentclass[10pt]{beamer}
\usetheme{metropolis}
\usepackage{graphicx}
\usepackage{listings}
\usepackage{booktabs}
\usepackage{xcolor}
\usepackage{hyperref}

\definecolor{codebg}{HTML}{F5F5F5}
\definecolor{codegreen}{HTML}{2E7D32}
\definecolor{codegray}{HTML}{757575}
\definecolor{codeblue}{HTML}{1565C0}

\lstdefinestyle{rstyle}{
  language=R, backgroundcolor=\color{codebg},
  basicstyle=\ttfamily\scriptsize,
  keywordstyle=\color{codeblue}\bfseries,
  stringstyle=\color{codegreen},
  commentstyle=\color{codegray}\itshape,
  breaklines=true, frame=single,
  rulecolor=\color{codegray!30}, showstringspaces=false, tabsize=2,
  morekeywords={library,install.packages,read_csv,glimpse,head,str,summary,
    plot,hist,barplot,boxplot,par,mfrow,col,main,xlab,ylab,pch,cex,lwd,
    abline,legend,text,lines,points,title,axis,data.frame,c,seq,rep,
    factor,as.Date,class,typeof,is.numeric,pipe,filter,mutate,select,
    group_by,summarise,arrange,tibble,ggplot,aes,geom_point,theme_minimal}
}
\lstdefinestyle{bashstyle}{
  language=bash, backgroundcolor=\color{codebg},
  basicstyle=\ttfamily\scriptsize,
  keywordstyle=\color{codeblue}\bfseries,
  commentstyle=\color{codegray}\itshape,
  breaklines=true, frame=single,
  rulecolor=\color{codegray!30}, showstringspaces=false,
}

\title{R Environment Setup \& Base Graphics}
\subtitle{Workshop 1 --- Module 8}
\author{Prof.Asoc.\ Endri Raco}
\institute{Data Visualization for Data Scientists \\ UNYT Tirana}
\date{2025/26}

\begin{document}

\begin{frame}
  \titlepage
\end{frame}

% ================================================================
\begin{frame}{Module Roadmap}
  \begin{enumerate}
    \item Installing R and RStudio
    \item The tidyverse ecosystem
    \item R data types and structures for visualization
    \item Base R graphics: \texttt{plot()}, \texttt{hist()}, \texttt{barplot()}, \texttt{boxplot()}
    \item The \texttt{par()} system and multi-panel layouts
    \item Base R vs ggplot2: when to use each
  \end{enumerate}
  \begin{exampleblock}{Learning Outcome}
    You will have a functioning R environment and be able to produce
    all four fundamental chart types using base R graphics, preparing
    the ground for ggplot2 in Workshop 2.
  \end{exampleblock}
\end{frame}

% ================================================================
\section{R \& RStudio Setup}
% ================================================================

\begin{frame}[fragile]{Installation Checklist}
  \begin{columns}[T]
    \begin{column}{0.48\textwidth}
      \textbf{Step 1: Install R}
\begin{lstlisting}[style=bashstyle]
# macOS (Homebrew)
brew install r

# Ubuntu / Debian
sudo apt install r-base

# Windows
# Download from https://cran.r-project.org
\end{lstlisting}
      \vspace{4pt}
      \textbf{Step 2: Install RStudio}\\
      {\small Download from \texttt{posit.co/download/rstudio-desktop}}
    \end{column}
    \begin{column}{0.48\textwidth}
      \textbf{Step 3: Install core packages}
\begin{lstlisting}[style=rstyle]
# One command installs 8+ packages
install.packages("tidyverse")

# Additional viz packages
install.packages(c(
  "viridis",
  "RColorBrewer",
  "ggthemes",
  "patchwork",
  "ggrepel",
  "scales"
))
\end{lstlisting}
    \end{column}
  \end{columns}
\end{frame}

% ================================================================
\begin{frame}{The RStudio IDE}
  \begin{center}
    \includegraphics[width=0.88\textwidth]{figures_m08/rstudio_layout}
  \end{center}
\end{frame}

% ================================================================
\begin{frame}[fragile]{RStudio Keyboard Shortcuts}
  \centering
  \begin{tabular}{lll}
    \toprule
    \textbf{Action} & \textbf{macOS} & \textbf{Windows} \\
    \midrule
    Run current line     & Cmd+Enter       & Ctrl+Enter \\
    Run entire script    & Cmd+Shift+Enter & Ctrl+Shift+Enter \\
    Insert pipe \texttt{|>} & Cmd+Shift+M  & Ctrl+Shift+M \\
    Insert assignment \texttt{<-} & Option+- & Alt+- \\
    Comment/uncomment    & Cmd+Shift+C     & Ctrl+Shift+C \\
    New R script         & Cmd+Shift+N     & Ctrl+Shift+N \\
    \bottomrule
  \end{tabular}

  \vspace{6pt}
  \begin{alertblock}{Pro-Tip}
    Use \texttt{Cmd+Shift+M} / \texttt{Ctrl+Shift+M} to insert the pipe
    operator instantly. This single shortcut will save hours over a term.
  \end{alertblock}
\end{frame}

% ================================================================
\section{The tidyverse Ecosystem}
% ================================================================

\begin{frame}{tidyverse: One \texttt{library()} Call, Eight Packages}
  \begin{center}
    \includegraphics[width=0.88\textwidth]{figures_m08/tidyverse_ecosystem}
  \end{center}
\end{frame}

% ================================================================
\begin{frame}[fragile]{The Pipe: Left-to-Right Data Flow}
  \begin{center}
    \includegraphics[width=0.95\textwidth]{figures_m08/pipe_workflow}
  \end{center}
  \vspace{2pt}
\begin{lstlisting}[style=rstyle]
# Traditional (nested, reads inside-out):
summarise(group_by(filter(data, Type == "Free"), Category), n = n())

# Pipe (linear, reads left-to-right):
data |>
  filter(Type == "Free") |>
  group_by(Category) |>
  summarise(n = n())
\end{lstlisting}
\end{frame}

% ================================================================
\section{R Data Types}
% ================================================================

\begin{frame}{Data Types for Visualization}
  \begin{center}
    \includegraphics[width=0.82\textwidth]{figures_m08/data_types}
  \end{center}
  \begin{alertblock}{Pro-Tip}
    Use \texttt{class(x)} to check a variable's type, and
    \texttt{str(df)} to see all column types in a data frame.
    Type mismatches (e.g., numbers stored as characters) are the most
    common source of plotting errors.
  \end{alertblock}
\end{frame}

% ================================================================
\begin{frame}[fragile]{Vectors, Data Frames, Tibbles}
\begin{lstlisting}[style=rstyle]
# Vector: 1-dimensional, single type
x <- c(10, 20, 30, 40, 50)
class(x)  # "numeric"

# Data frame: 2D table (base R)
df <- data.frame(
  name = c("Alice", "Bob", "Carol"),
  score = c(88, 92, 75),
  grade = factor(c("A", "A", "B"))
)

# Tibble: modern data frame (tidyverse)
library(tibble)
tb <- tibble(
  name  = c("Alice", "Bob", "Carol"),
  score = c(88, 92, 75),
  grade = factor(c("A", "A", "B"))
)

# Key difference: tibbles print compactly
# and never convert strings to factors
print(tb)
\end{lstlisting}
\end{frame}

% ================================================================
\begin{frame}[fragile]{Loading and Inspecting Data}
\begin{lstlisting}[style=rstyle]
library(readr)

# Read CSV (tidyverse — fast, type-guessing)
apps <- read_csv("googleplaystore.csv")

# Quick inspection
dim(apps)        # rows x columns
glimpse(apps)    # compact column view
head(apps, 5)    # first 5 rows
summary(apps)    # five-number summaries

# Check for problems
problems(apps)   # parsing issues
sum(is.na(apps)) # total missing values

# Useful for viz planning:
# - How many rows? (affects point density)
# - Column types? (determines chart choice)
# - Missing values? (need filtering/imputation)
\end{lstlisting}
\end{frame}

% ================================================================
\section{Base R Graphics}
% ================================================================

\begin{frame}{The Four Fundamental Chart Types}
  \begin{center}
    \includegraphics[width=0.88\textwidth]{figures_m08/par_system}
  \end{center}
  \begin{exampleblock}{Data Insight}
    Base R produces functional charts instantly with zero package
    dependencies. These four functions --- \texttt{plot()}, \texttt{hist()},
    \texttt{barplot()}, \texttt{boxplot()} --- cover 80\% of exploratory
    visualization needs.
  \end{exampleblock}
\end{frame}

% ================================================================
\begin{frame}[fragile]{Scatterplot: \texttt{plot(x, y)}}
  \begin{columns}[T]
    \begin{column}{0.52\textwidth}
\begin{lstlisting}[style=rstyle]
# Minimal call
plot(mpg$displ, mpg$hwy)

# With customisation
plot(mpg$displ, mpg$hwy,
  main = "Engine vs Fuel Economy",
  xlab = "Displacement (L)",
  ylab = "Highway MPG",
  pch  = 16,       # solid circle
  col  = "steelblue",
  cex  = 0.8)      # point size

# Add regression line
abline(
  lm(hwy ~ displ, data = mpg),
  col = "red", lwd = 2)

# Add legend
legend("topright",
  legend = "OLS fit",
  col = "red", lwd = 2)
\end{lstlisting}
    \end{column}
    \begin{column}{0.45\textwidth}
      \includegraphics[width=\textwidth]{figures_m08/base_scatter}

      \vspace{4pt}
      {\small Key parameters:
      \texttt{pch} = point shape (16 = filled),
      \texttt{col} = colour,
      \texttt{cex} = size multiplier,
      \texttt{lwd} = line width.}
    \end{column}
  \end{columns}
\end{frame}

% ================================================================
\begin{frame}[fragile]{Histogram: \texttt{hist(x)}}
  \begin{columns}[T]
    \begin{column}{0.52\textwidth}
\begin{lstlisting}[style=rstyle]
# Minimal call
hist(apps$Rating)

# With customisation
hist(apps$Rating,
  breaks = 30,          # bin count
  col    = "lightblue",
  border = "white",
  main   = "App Rating Distribution",
  xlab   = "Rating (1-5)",
  ylab   = "Frequency",
  las    = 1)           # horizontal labels

# Add density curve overlay
lines(density(apps$Rating,
  na.rm = TRUE),
  col = "red", lwd = 2)

# Add vertical mean line
abline(v = mean(apps$Rating,
  na.rm = TRUE),
  col = "blue", lwd = 2,
  lty = 2)  # dashed
\end{lstlisting}
    \end{column}
    \begin{column}{0.45\textwidth}
      \includegraphics[width=\textwidth]{figures_m08/base_hist}

      \vspace{4pt}
      {\small Key parameters:
      \texttt{breaks} = number of bins,
      \texttt{border} = bar edge colour,
      \texttt{las = 1} = horizontal tick
      labels.}
    \end{column}
  \end{columns}
\end{frame}

% ================================================================
\begin{frame}[fragile]{Bar Chart: \texttt{barplot()}}
  \begin{columns}[T]
    \begin{column}{0.52\textwidth}
\begin{lstlisting}[style=rstyle]
# Count categories
counts <- table(apps$Type)

# Minimal bar chart
barplot(counts)

# Customised horizontal bar
barplot(sort(counts),
  horiz   = TRUE,
  col     = "steelblue",
  border  = NA,
  main    = "Apps by Type",
  xlab    = "Count",
  las     = 1,
  cex.names = 0.8)

# Grouped bar chart
barplot(
  matrix(c(30, 50, 20, 40), nrow = 2),
  beside = TRUE,
  col = c("steelblue", "coral"),
  names.arg = c("Cat A", "Cat B"),
  legend.text = c("Free", "Paid"))
\end{lstlisting}
    \end{column}
    \begin{column}{0.45\textwidth}
      \includegraphics[width=\textwidth]{figures_m08/base_bar}

      \vspace{4pt}
      {\small Key parameters:
      \texttt{horiz = TRUE} flips to horizontal,
      \texttt{beside = TRUE} groups bars
      side by side (vs.\ stacked),
      \texttt{border = NA} removes bar edges.}
    \end{column}
  \end{columns}
\end{frame}

% ================================================================
\begin{frame}[fragile]{Box Plot: \texttt{boxplot()}}
  \begin{columns}[T]
    \begin{column}{0.52\textwidth}
\begin{lstlisting}[style=rstyle]
# Formula syntax: y ~ group
boxplot(Rating ~ Type, data = apps,
  main   = "Rating by App Type",
  ylab   = "Rating",
  col    = c("lightblue", "salmon"),
  border = "grey40",
  notch  = TRUE,   # CI for median
  las    = 1)

# Add mean points
means <- tapply(apps$Rating,
  apps$Type, mean, na.rm = TRUE)
points(1:2, means,
  pch = 18, col = "red", cex = 1.5)

# Horizontal boxplot
boxplot(Rating ~ Type, data = apps,
  horizontal = TRUE)
\end{lstlisting}
    \end{column}
    \begin{column}{0.45\textwidth}
      \includegraphics[width=\textwidth]{figures_m08/base_box}

      \vspace{4pt}
      {\small The formula \texttt{y \textasciitilde{} group}
      is R's standard interface for
      conditioning: ``y split by group.''
      Used in boxplot, lm, aov, etc.}
    \end{column}
  \end{columns}
\end{frame}

% ================================================================
\section{The par() System}
% ================================================================

\begin{frame}[fragile]{Multi-Panel Layouts with \texttt{par(mfrow)}}
\begin{lstlisting}[style=rstyle]
# Set up 2x2 grid (fills row by row)
par(mfrow = c(2, 2))

# Panel 1: scatter
plot(mpg$displ, mpg$hwy, pch = 16, cex = 0.6,
     main = "Scatter", xlab = "Displacement", ylab = "MPG")

# Panel 2: histogram
hist(mpg$hwy, breaks = 20, col = "lightblue", border = "white",
     main = "Distribution", xlab = "Highway MPG")

# Panel 3: bar chart
barplot(table(mpg$class), col = "steelblue", border = NA,
        main = "Vehicle Class", las = 2, cex.names = 0.7)

# Panel 4: boxplot
boxplot(hwy ~ cyl, data = mpg, col = "lightblue",
        main = "MPG by Cylinders", xlab = "Cylinders", ylab = "MPG")

# Reset to single plot
par(mfrow = c(1, 1))
\end{lstlisting}
\end{frame}

% ================================================================
\begin{frame}[fragile]{Customising \texttt{par()} Parameters}
  \centering
  \begin{tabular}{lp{5.5cm}}
    \toprule
    \textbf{Parameter} & \textbf{Effect} \\
    \midrule
    \texttt{mfrow = c(r, c)} & Grid with r rows, c columns \\
    \texttt{mar = c(b, l, t, r)} & Margins in lines (bottom, left, top, right) \\
    \texttt{mgp = c(3, 1, 0)} & Axis title, labels, line positions \\
    \texttt{cex.main = 1.2} & Title size multiplier \\
    \texttt{cex.axis = 0.8} & Tick label size multiplier \\
    \texttt{bg = "white"} & Background colour \\
    \texttt{las = 1} & Horizontal tick labels \\
    \bottomrule
  \end{tabular}

  \vspace{6pt}
  \begin{alertblock}{Pro-Tip}
    Always save and restore \texttt{par()}: \\
    \texttt{old\_par <- par(mfrow = c(2,2))} \\
    \texttt{\# ... plots ...} \\
    \texttt{par(old\_par)}
  \end{alertblock}
\end{frame}

% ================================================================
\section{Base R vs ggplot2}
% ================================================================

\begin{frame}{When to Use Each}
  \begin{center}
    \includegraphics[width=0.92\textwidth]{figures_m08/base_vs_ggplot}
  \end{center}
\end{frame}

% ================================================================
\begin{frame}{Base R vs ggplot2: Decision Guide}
  \centering
  \begin{tabular}{lll}
    \toprule
    \textbf{Criterion} & \textbf{Base R} & \textbf{ggplot2} \\
    \midrule
    Speed to first plot & Fastest (1 line) & Slightly more setup \\
    Customisation      & Manual, verbose    & Declarative, layered \\
    Faceting           & \texttt{par(mfrow)} & \texttt{facet\_wrap()} \\
    Themes             & Manual \texttt{par()} & \texttt{theme\_*()} library \\
    Colour scales      & Manual vectors     & \texttt{scale\_*} functions \\
    Publication quality & Requires effort   & Near-ready by default \\
    Learning curve     & Shallow start      & Steeper, but compounds \\
    \bottomrule
  \end{tabular}

  \vspace{6pt}
  \begin{exampleblock}{Data Insight}
    Use \textbf{base R} for quick exploratory checks (``what does this
    variable look like?''). Use \textbf{ggplot2} for anything that will
    be seen by others. We cover ggplot2 in depth in Workshop 2.
  \end{exampleblock}
\end{frame}

% ================================================================
\section{Synthesis}
% ================================================================

\begin{frame}{Key Takeaways}
  \begin{enumerate}
    \item \textbf{RStudio} is the standard IDE: learn the four panes and
          the pipe shortcut (\texttt{Cmd/Ctrl+Shift+M}).
    \item \textbf{tidyverse} loads eight essential packages with one
          \texttt{library()} call. Master \texttt{readr} + \texttt{dplyr}
          + \texttt{ggplot2} first.
    \item \textbf{Data types} determine chart choices: check with
          \texttt{class()} and \texttt{str()} before plotting.
    \item \textbf{Base R graphics} (\texttt{plot}, \texttt{hist},
          \texttt{barplot}, \texttt{boxplot}) produce instant exploratory
          charts with zero dependencies.
    \item \textbf{par(mfrow)} creates multi-panel layouts; always save
          and restore \texttt{par()} settings.
    \item Base R is for \textbf{exploration}; ggplot2 (Workshop 2) is for
          \textbf{explanation and publication}.
  \end{enumerate}
\end{frame}

% ================================================================
\begin{frame}{Essential Readings for This Module}
  \begin{itemize}
    \item Wickham, H. \& Grolemund, G. (2023). \emph{R for Data Science},
          2nd ed., Chapters 1--3 (Setup, Workflow, Data Transformation).
          \texttt{https://r4ds.hadley.nz}
    \item Murrell, P. (2018). \emph{R Graphics}, 3rd ed., Chapters 1--3
          (Base Graphics System).
    \item Peng, R.~D. (2016). \emph{Exploratory Data Analysis with R},
          Chapters 4--6 (Plotting Systems).
  \end{itemize}

  \vspace{6pt}
  \textbf{Next Module}: Python Environment Setup \& Matplotlib Basics (W01-M09)
\end{frame}

\end{document}
